\section{Introduction}
\label{sec:introduction}

Retrieval-augmented generation (RAG) grounds model outputs in external
evidence~\cite{lewis2020retrieval}, while multimodal extensions retrieve images
and text for knowledge-intensive visual reasoning
~\cite{chen2022murag,yasunaga2023retrieval,long2024gemkr}. However, the benefit
of evidence-grounded generation is fundamentally constrained by candidate
coverage. Once answer-supporting evidence is absent from the retrieved pool,
downstream reranking or selection can only reorder inadequate candidates. A
reliable system must therefore decide not only which evidence to use, but also
whether the current pool is sufficient, which partial evidence should persist,
and when retrieval should be redirected.

This problem is acute in multimodal endoscopy. Different pathologies or
anatomical sites can appear similar, while question text may provide weak
retrieval cues~\cite{li2023llavamed,liu2025endobench}. Fixed fusion can
overweight an uninformative modality, and unrestricted retrieval can return
visually plausible but anatomically incompatible evidence. Rewriting may
improve coverage, but a failed rewrite can shift retrieval to another poor
neighborhood. Reranking cannot recover absent evidence, while discarding prior
failures removes clues for avoiding the same direction. Such reliability
matters in medical RAG, where technical development still outpaces systematic
clinical validation~\cite{yang2026medicalrag}.

Figure~\ref{fig:motivation} motivates the distinction between myopic and
stateful retrieval control. A memoryless controller reformulates only
from the current round and may revisit a previously explored, misleading
retrieval neighborhood. A stateful controller can instead retain useful
partial evidence and failed retrieval traces. The former can accumulate as
answer-supporting context; the latter can inform subsequent reformulation
without contaminating the final generation context.

\begin{figure*}[t]
  \centering
  \includegraphics[width=\textwidth]{fig1_motivation_submission.png}
  \caption{Design motivation for myopic versus stateful retrieval control. A memoryless
  controller can repeatedly explore the same misleading neighborhood.
  AgenticRL maintains retained evidence $M_t^+$ as generation-visible context
  and failed retrieval traces $M_t^-$ as controller-only history for subsequent
  reformulation.}
  \label{fig:motivation}
\end{figure*}

Recent work studies RL for multi-round search, retriever-aware rewriting, and
multimodal RAG~\cite{jin2025searchr1,song2025r1searcher,li2026ttr,
lin2026craf,zhao2026mmragrft}. These controllers reformulate from the current
round only, so what a round has established is lost once the next query is
issued. We instead train one smaller policy to manage evidence persistence and
further search under a fixed budget over a frozen multimodal stack. Failed
traces remain available for control but hidden from answer generation.

We present \textbf{AgenticRL}, a closed-loop framework for controlling this
frozen environment. The retrieval front end uses anatomical semantic
partitioning and optionally applies query-dependent modality weighting before
retrieving multimodal candidates. A base controller selects evidence and
chooses \textsc{Accept} or \textsc{Rewrite}. The complete inference strategy
wraps this controller in a fixed three-round refinement process. Its
generation-visible state $M_t^+$ stores the latest non-empty retained set,
whereas previous queries and failed candidate traces form the controller-only
state $M_t^-$. An initial \textsc{Rewrite} deterministically initializes
$M^+$ with the three highest-ranked candidates; subsequent non-empty
\textsc{Keep} selections replace the current best evidence state.

Both actions use a shared generator-derived utility. \textsc{Accept} scores
retained evidence; \textsc{Rewrite} first executes the complete frozen
retrieval path and then scores its returned evidence. This penalizes fluent
rewrites that fail against the deployed index. An SFT--RFT--GRPO curriculum
optimizes only the controller.

We evaluate on EndoBench~\cite{liu2025endobench}, which contains $6{,}832$ VQA
pairs across four endoscopic scenarios and 12 clinical tasks. Under a shared
frozen stack, we compare one-shot RAG, single-round control, memoryless
iterative control, and memory-aware control, and we isolate the state itself
by a protocol-matched ablation and by separating questions on which memory
becomes active from those on which it never does. The results show that within
budgeted execution it is cross-round state, not iteration, that makes iterative
retrieval pay off: memoryless iteration is no better than a single round
(39.56\% vs.\ 40.12\%), whereas AgenticRL reaches 42.84\%, i.e.\ $+3.28$
points over memoryless iterative control (exact McNemar $p<10^{-12}$) and
$+2.72$ over single-round control, with the gain concentrated exactly where the
state is used. A 4B controller with a frozen 8B generator thereby surpasses the
41.69\% GPT-4o accuracy that EndoBench reports by 1.15 points, and also exceeds
its 34B and 80B medical-specialist models; those systems answer without
retrieval, so the comparison is contextual, but it shows that state-carrying
retrieval control reaches this accuracy band more cheaply than model scale.

Our contributions are:
\begin{itemize}
  \item We formulate frozen-stack multimodal RAG as budgeted closed-loop
  control, where one policy jointly selects candidate evidence, updates a
  current best evidence state, and redirects the deployed retriever when
  information remains missing.
  \item We introduce trajectory-aware retrieval memory that separates
  generation-visible retained evidence $M^+$ from controller-only failed
  retrieval traces $M^-$, together with the deterministic budgeted wrapper that
  makes retained evidence reusable across rounds. This lets the latest useful
  evidence persist while refuted retrieval directions inform later
  reformulation without ever entering the answer context.
  \item We place heterogeneous \textsc{Accept} and \textsc{Rewrite} actions on
  a shared generator-derived scale after executing their actual retrieval
  consequences, and optimize the controller with an SFT--RFT--GRPO curriculum
  without updating the retrieval environment or generator.
  \item On the full EndoBench benchmark we show that, under a fixed budget,
  cross-round state rather than iteration is what makes repeated retrieval pay
  off: memoryless iterative control does not improve over a single round,
  whereas the complete evidence-preserving strategy adds $3.28$ points over it
  ($p<10^{-12}$) and $2.72$ points over single-round control, exceeds the
  benchmark-reported GPT-4o accuracy with a 4B controller, and produces no
  change on the questions where the state never activates. Matched diagnostics
  further delimit how much of this gain the state itself carries.
\end{itemize}
